% Section 5 -- Runtime Execution

\subsection{From Descriptors to Runtime State}

Phase 3 stopped at descriptor emission: the compiler and validator
produced \texttt{PeerHandoffDescriptor}s, but the runtime still served
every admit from host memory. The final runtime keeps the compiler and
validator unchanged and only changes how those descriptors are
consumed.

At load time, the projected descriptors are handed to
\texttt{MemPartitionBufferStorage}, indexed by
$(\texttt{dst\_lane}, \texttt{round\_idx}, \texttt{partition\_id})$,
and checked for duplicates. Runtime initialization then attempts to
enable CUDA peer access in both directions for every active device
pair. Two environment controls gate the feature:
\texttt{GEGE\_STATEFLOW\_PEER\_RUNTIME=\{off,auto,on\}} controls whether
peer execution is disabled, opportunistic, or required, and
\texttt{GEGE\_STATEFLOW\_PEER\_RUNTIME\_SCOPE=\{all,embeddings\}}
controls whether the peer path is applied to both embeddings and
optimizer state or to embeddings only.

\subsection{In-Place Peer Execution}

Our first Phase 4 prototype allocated a full staged destination buffer
per device. That design was simple but too memory-hungry for
\SI{24}{GB} boards. The final Phase 4b runtime instead performs
\emph{in-place, slot-aligned} peer admits.

At each round boundary, each lane materializes only its own next-state
layout locally, while the runtime consults a precomputed per-round
``globally required'' set to decide which evictions truly leave the
next frontier. This removes the old cross-device next-state barrier.
The source side snapshots only those partitions that will be exported
to other lanes into small per-partition source-scratch tensors and
publishes their slot locations plus a CUDA ready event. Destinations
wait only on the specific source lanes they consume, not on a global
all-lane barrier. This decouples peer reads from subsequent local slot
reuse without allocating another full device-sized buffer.

The destination side then realizes the planner's next-state slot layout
exactly. Retained partitions are permuted into their target slots using
free evict slots as temporary scratch when necessary; partitions that
leave the global next-state frontier are written back to host; and each
new admit is filled either by \texttt{cudaMemcpyPeerAsync} from the
source scratch or by the legacy host path. Because the runtime
materializes the planner's slot layout rather than reusing slots
opportunistically, the compiler's \texttt{src\_slot} and
\texttt{dst\_slot} remain valid at execution time.

\subsection{Safety and Telemetry}

The validator still enforces the structural invariants of
Section~\ref{sec:ir:validator}; the runtime adds dynamic checks that the
projected source lane, destination lane, and slot IDs match the actual
swap-time state. Any mismatch falls back to host reload and increments
a visible counter rather than risking a corrupt swap. This made the
bring-up path auditable: we could distinguish ``descriptor emitted but
not executed,'' ``descriptor executed with fallback,'' and ``descriptor
executed exactly.''

Each epoch reports
\texttt{peer\_bytes\_executed}, \texttt{host\_fallback\_bytes},
\texttt{peer\_copy\_count}, \texttt{descriptor\_mismatch\_count}, and
\texttt{peer\_sync\_wait\_ms}, with per-device breakdowns. These are
the counters used in Section~\ref{sec:eval} to show that the final
Phase 4b runtime executes the planned peer relays with zero fallback
and zero slot mismatch.
